\documentclass[11pt, a4paper]{article}

% ── PACKAGES ──────────────────────────────────────────────────
\usepackage[utf8]{inputenc}
\usepackage[T1]{fontenc}
\usepackage{geometry}
\geometry{margin=2.5cm}
\usepackage{hyperref}
\hypersetup{
    colorlinks=true,
    linkcolor=blue,
    urlcolor=blue,
    citecolor=blue,
    pdftitle={FOXA1/EZH2 IHC Validation
              Partnership Proposal v3.0},
    pdfauthor={Eric Robert Lawson / OrganismCore}
}
\usepackage{booktabs}
\usepackage{array}
\usepackage{parskip}
\usepackage{microtype}
\usepackage{authblk}
\usepackage{titlesec}
\usepackage{fancyhdr}
\usepackage{xcolor}
\usepackage{enumitem}
\usepackage{mdframed}
\usepackage{multirow}
\usepackage{amsmath}

% ── COLOURS ───────────────────────────────────────────────────
\definecolor{headerblue}{RGB}{20,60,140}
\definecolor{findingbox}{RGB}{240,250,240}
\definecolor{askbox}{RGB}{235,245,255}
\definecolor{equitybox}{RGB}{255,245,230}
\definecolor{statsbox}{RGB}{248,240,255}
\definecolor{honestybox}{RGB}{255,248,220}
\definecolor{dynamicbox}{RGB}{235,255,245}
\definecolor{tmabox}{RGB}{240,248,255}
\definecolor{infobox}{RGB}{235,245,255}

% ── PAGE STYLE ────────────────────────────────────────────────
\pagestyle{fancy}
\fancyhf{}
\lhead{\footnotesize FOXA1/EZH2 IHC Validation
       Partnership Proposal v3.0}
\rhead{\footnotesize OrganismCore \textbar{}
       2026-03-06}
\rfoot{\small Page \thepage}
\renewcommand{\headrulewidth}{0.4pt}
\setlength{\headheight}{24pt}
\addtolength{\topmargin}{-10pt}

% ── SECTION FORMATTING ────────────────────────────────────────
\titleformat{\section}
  {\large\bfseries\color{headerblue}}
  {\thesection}{1em}{}[\titlerule]
\titleformat{\subsection}
  {\normalsize\bfseries}{\thesubsection}{1em}{}

% ══════════════════════════════════════════════════════════════
\title{
    \vspace{-1.5em}
    {\color{headerblue}\textbf{Validation
    Partnership Proposal --- Version 3.0
    (Final)}}\\[0.4em]
    \large\textbf{FOXA1/EZH2 Dual IHC as a
    Universal Breast Cancer Treatment
    Classifier}\\[0.3em]
    \normalsize Two stains. One ratio.
    Six subtypes. Same-day results.
    \$50--\$100 per patient.
    TMA-compatible.\\[0.5em]
    \small\textit{A retrospective concordance
    study against PAM50 in archival FFPE tissue
    --- the single remaining step to clinical
    deployment. Evidence base, honest limitations,
    and TMA compatibility fully stated.}
}

\author[1]{Eric Robert Lawson}
\affil[1]{OrganismCore (Independent Research)
    \textbar{}
    \href{mailto:OrganismCore@proton.me}
    {OrganismCore@proton.me} \textbar{}
    ORCID:
    \href{https://orcid.org/0009-0002-0414-6544}
    {0009-0002-0414-6544}}

\date{March 6, 2026}

% ══════════════════════════════════════════════════════════════
\begin{document}
\maketitle
\thispagestyle{fancy}

\begin{mdframed}[backgroundcolor=infobox,
                 linecolor=blue!40!black,
                 linewidth=1pt]
\textbf{VERSION HISTORY:}
v1.0 (2026-03-06): Initial release.
v2.0 (2026-03-06): Dynamic range rationale;
honest limitations section added.
v3.0 (2026-03-06): TMA compatibility section
added. Final version for distribution.
\end{mdframed}

\vspace{0.5em}
\noindent\rule{\textwidth}{1.5pt}

% ── SUMMARY BOX ───────────────────────────────────────────────
\begin{mdframed}[backgroundcolor=findingbox,
                 linecolor=green!50!black,
                 linewidth=1.5pt]
\textbf{In brief:}

The ratio FOXA1~H-score $\div$ EZH2~H-score,
computed from two standard IHC stains, was
derived from first principles using Waddington
attractor geometry applied to 19,542 single
cancer cells. It correctly orders all six major
breast cancer subtypes and maps each to its
therapeutic treatment logic. Computationally
validated across ${\sim}7{,}500$ patients in
seven independent datasets on four measurement
platforms. Zero biological contradictions.

\medskip
\textbf{The computational case is complete.}

\medskip
The test costs \$50--\$100 per patient. It
runs on equipment present in any pathology
laboratory worldwide. It produces results
the same day as biopsy processing. It requires
no RNA, no proprietary platform, no reference
laboratory. It is \textbf{compatible with
tissue microarray (TMA) format} --- the
standard archival tissue infrastructure at
most major cancer centres.

\medskip
One study is needed before clinical deployment:
a retrospective IHC concordance study against
PAM50 in $n=300$--$400$ archived FFPE cases
(whole section or TMA). No new patients.
No treatment changes. Estimated cost:
\$5,000--\$15,000 in consumables.
Timeline: 3--6 months.

\medskip
\textbf{We are seeking one or more pathology
or oncology partners to conduct this study.}
\end{mdframed}

\vspace{0.8em}

% ══════════════════════════════════════════════════════════════
\section{The Problem}
% ══════════════════════════════════════════════════════════════

Breast cancer causes ${\sim}685{,}000$ deaths
per year globally. The majority occur where
PAM50/Prosigna is unavailable.

\begin{center}
\begin{tabular}{p{4.2cm}p{4.5cm}p{4.5cm}}
\toprule
& \textbf{PAM50/Prosigna} &
\textbf{FOXA1/EZH2 IHC ratio} \\
\midrule
Cost per patient &
\$3,000--\$4,000 &
\$50--\$100 \\
Platform &
NanoString nCounter
(\$150k--\$250k capital) &
Standard IHC \\
RNA required &
Yes (degrades; strict handling) &
No \\
Laboratory type &
Specialised, certified &
Any pathology laboratory \\
Turnaround &
Days to weeks &
Same day \\
TMA compatible &
No (requires fresh extraction) &
Yes --- fully compatible \\
Output &
Subtype label only &
Subtype $+$ mechanism $+$
treatment logic \\
Global access &
Major centres in
wealthy countries only &
Universal \\
\bottomrule
\end{tabular}
\end{center}

Without molecular subtyping, ER$+$ patients
(LumA and LumB --- different mechanisms,
different optimal treatment) receive identical
treatment. Triple-negative patients (TNBC and
CL --- opposite treatment implications) receive
identical treatment. The ratio resolves both
ambiguities from two stains on a slide already
being prepared.

% ══════════════════════════════════════════════════════════════
\section{The Finding and Its Derivation}
% ══════════════════════════════════════════════════════════════

\subsection{Biological rationale}

FOXA1 is the pioneer transcription factor
that opens chromatin for the luminal identity
programme. EZH2 (PRC2 catalytic subunit)
silences that programme by depositing H3K27me3
marks on the FOXA1, GATA3, and ESR1 promoters.
Their ratio directly measures the balance
between luminal identity and its epigenetic
suppression --- the biological axis determining
both subtype and therapeutic vulnerability.

The ratio was derived from Waddington landscape
attractor geometry applied to 19,542 single
cancer cells (GSE176078, Wu et al.\ 2021,
\textit{Nature Genetics}), with all predictions
locked before any confirmatory analysis or
literature review. The predicted ordering
was confirmed exactly. Zero biological
contradictions were found across all seven
datasets.

\subsection{Validated subtype ordering and
treatment logic}

\begin{center}
\begin{tabular}{p{2.2cm}p{1.7cm}p{8.3cm}}
\toprule
\textbf{Subtype} &
\textbf{scRNA-seq ratio$^{\dagger}$} &
\textbf{Treatment logic} \\
\midrule
Luminal A & $9.38$ &
Luminal programme intact.
CDK4/6 inhibitor $+$ ET. \\[0.3em]
Luminal B & $8.10$ &
Chromatin lock.
HDACi (entinostat) $+$ ET. \\[0.3em]
HER2-enriched & $3.34$ &
HER2 amplicon.
Anti-HER2 first, ET thereafter. \\[0.3em]
TNBC & $0.52$ &
EZH2 epigenetic lock ($+189\%$).
Tazemetostat $\rightarrow$ fulvestrant. \\[0.3em]
Claudin-low & $0.10$ &
No lock to dissolve.
Immune only.
Anti-TIGIT $\rightarrow$ anti-PD-1. \\[0.3em]
ILC (exception) &
Inverted &
FOXA1 hyperactivated above LumA.
Structural lock. Fulvestrant $>$ AI. \\
\bottomrule
\end{tabular}
\end{center}

\small $^{\dagger}$\textit{scRNA-seq units
only. Not IHC H-score values. IHC cut-points
to be determined in the calibration study.}
\normalsize

% ══════════════════════════════════════════════════════════════
\section{The Computational Evidence Base}
% ══════════════════════════════════════════════════════════════

\begin{mdframed}[backgroundcolor=statsbox,
                 linecolor=purple!40!black,
                 linewidth=1pt]
\begin{itemize}[noitemsep]
    \item \textbf{Subtype ordering:}
    TCGA-BRCA $n=837$,
    KW $p=2.87\times10^{-103}$;
    METABRIC LumA vs LumB
    $p=8.47\times10^{-67}$, $n=1{,}980$
    \item \textbf{Survival:} TCGA KM
    $p=0.0031$ ($n=1{,}218$); METABRIC
    KM $p<0.0001$ ($n=1{,}980$);
    GSE96058 KM $p\approx0$ ($n=3{,}273$,
    336 events, 52 months follow-up).
    Three independent cohorts. Two platforms.
    \item \textbf{ROC classification:}
    LumA vs Basal AUC $0.828$--$0.901$;
    LumA vs LumB AUC $0.796$--$0.873$
    (replicated in METABRIC and TCGA
    independently)
    \item \textbf{Protein-level confirmation:}
    CPTAC-BRCA iTRAQ MS, $n=122$;
    correct subtype ordering at protein level;
    Spearman $r=-0.492$, $p<0.0001$
    \item \textbf{Treatment response:}
    GSE25066 chemotherapy $n=508$,
    $p<0.0001$; METABRIC endocrine therapy
    $n=1{,}104$, $p=0.0027$,
    $\Delta=18.7$ months RFS
    \item \textbf{7 independent datasets.
    4 measurement platforms.
    0 biological contradictions.}
\end{itemize}
\end{mdframed}

% ══════════════════════════════════════════════════════════════
\section{Why IHC Is Expected to Succeed:
The Dynamic Range Argument}
% ══════════════════════════════════════════════════════════════

\begin{mdframed}[backgroundcolor=dynamicbox,
                 linecolor=green!60!black,
                 linewidth=1pt]
The protein ordering is confirmed by CPTAC
mass spectrometry. The MS kappa analysis for
subtype classification failed --- not because
the biology is wrong, but because iTRAQ MS
compressed all breast cancer subtypes into
\textbf{0.38 log$_2$ units} of total range.
At this scale, noise exceeds signal and
threshold classification is impossible.

IHC H-scores span \textbf{0--300} ---
approximately \textbf{800$\times$} more
dynamic range. The protein signal that is
unresolvable at 0.38 MS units is expected
to be clearly separable at IHC scale.

\textbf{The MS kappa failure is a reason to
be more confident in IHC, not less.} The
signal exists. The instrument used to test
classification lacked the range. IHC has
the range.
\end{mdframed}

% ══════════════════════════════════════════════════════════════
\section{What Is Established and What Is Not}
\label{sec:honesty}
% ══════════════════════════════════════════════════════════════

\begin{mdframed}[backgroundcolor=honestybox,
                 linecolor=orange!60!black,
                 linewidth=1pt]
\textbf{We state this explicitly.}

\medskip
\textbf{What IS established:}
\begin{itemize}[noitemsep]
    \item Correct subtype ordering in RNA
    and protein across multiple cohorts
    \item Survival stratification in four
    independent analyses (${\sim}7{,}500$
    patients)
    \item Treatment response prediction in
    two independent treatment datasets
    \item Protein signal confirmed at
    protein level (CPTAC MS, $n=122$)
    \item Both antibodies established in
    routine pathology use globally
    \item Biological mechanism independently
    confirmed (Schade et al.\
    \textit{Nature} 2024; Toska et al.\
    \textit{Nature Medicine} 2017)
\end{itemize}

\medskip
\textbf{What is NOT yet established:}
\begin{itemize}[noitemsep]
    \item \textbf{IHC H-score cut-points.}
    RNA-space values (9.38, 8.10, etc.)
    are not IHC thresholds. No formula
    converts between them. Cut-points
    must be determined from scratch in
    the calibration study.
    \item \textbf{IHC concordance with
    PAM50.} AUC values come from RNA-level
    analysis. IHC performance is expected
    to be strong based on the dynamic
    range argument but has not been
    measured.
    \item \textbf{Inter-observer
    reliability.} Not yet assessed.
    \item \textbf{CL below TNBC in IHC.}
    Confirmed in scRNA-seq (cancer cells
    only: CL 0.10, TNBC 0.52). Not
    confirmed in bulk RNA (stromal artefact
    --- explained). IHC should replicate
    scRNA-seq if cancer cells are scored
    exclusively. Requires confirmation.
\end{itemize}

\medskip
\textbf{None of these are conceptual gaps.}
They are instrument calibration steps ---
the same steps applied to ER, PR, HER2,
and Ki67 before routine clinical adoption.
\end{mdframed}

% ══════════════════════════════════════════════════════════════
\section{Tissue Microarray Compatibility}
% ══════════════════════════════════════════════════════════════

\begin{mdframed}[backgroundcolor=tmabox,
                 linecolor=blue!50!black,
                 linewidth=1pt]
\textbf{TMA FORMAT IS FULLY COMPATIBLE WITH
THIS PROTOCOL.}

Many major cancer centres hold archival breast
cancer tissue in tissue microarray format.
This protocol applies without modification to
TMA cores. No format conversion is required.

\textbf{TMA-specific requirements:}
\begin{itemize}[noitemsep]
    \item Minimum core diameter: 1.0~mm;
    1.5--2.0~mm preferred
    \item Minimum two cores per case per
    stain; cases with one evaluable core
    flagged and reported separately
    \item Score invasive cancer cell nuclei
    only within each core
    \item Two-core averaging: report mean
    ratio as case value; flag cases where
    two core ratios differ by $>0.5$ log
    units (may indicate heterogeneity or
    stromal core placement)
    \item Cut 3+ sequential sections at
    run start to ensure paired FOXA1 and
    EZH2 sections for each core
\end{itemize}

\textbf{TMA advantage:} A TMA of 400 cases
requires a single pair of slides per
staining run. All cases are processed
under identical conditions simultaneously,
reducing inter-case variability. Pathologist
scoring time is substantially reduced.
For institutions with TMA infrastructure,
this is the recommended format for the
calibration study.

\textbf{CL note:} TMA cores from
mesenchymal-rich claudin-low tumours may
have low cancer cell density. Cores with
$<50$ evaluable invasive cancer cell
nuclei should be recorded as unevaluable.
The stromal contamination warning (Protocol
Specification v3.0, Section~6) applies
with equal force to TMA cores.
\end{mdframed}

% ═════════════════════════════════��════════════════════════════
\section{The Partnership Ask}
% ══════════════════════════════════════════════════════════════

\begin{mdframed}[backgroundcolor=askbox,
                 linecolor=blue!50!black,
                 linewidth=1.5pt]
\textbf{What we are asking for:}
\begin{enumerate}
    \item \textbf{Access to archived FFPE
    breast cancer tissue} (whole section
    or TMA format) from $n=300$--$400$
    cases with known PAM50 subtype
    classification
    \item \textbf{FOXA1 and EZH2 IHC} run
    on those sections using the attached
    Protocol Specification v3.0
    \item \textbf{H-score scoring} by one
    or two pathologists, blinded to PAM50
    subtype
    \item \textbf{Concordance analysis}
    comparing IHC ratio classification
    to PAM50, using Youden J optimisation
    to determine cut-points from scratch
\end{enumerate}

\textbf{What we provide:}
\begin{enumerate}
    \item Protocol Specification v3.0
    (attached) --- antibody clones,
    catalogue numbers, antigen retrieval,
    H-score method, TMA instructions,
    claudin-low scoring warning
    \item Statistical analysis framework
    (concordance, ROC, cut-point
    optimisation, training/test split,
    inter-observer kappa)
    \item Complete public computational
    validation dataset and scripts
    \item Co-authorship on the resulting
    publication
\end{enumerate}

\textbf{What this does NOT require:}
\begin{itemize}[noitemsep]
    \item No new patient recruitment
    \item No treatment changes
    \item No new laboratory equipment
    \item No RNA extraction or
    bioinformatics
    \item Minimal regulatory burden
    (retrospective tissue research,
    minimal-risk pathway in most
    jurisdictions)
\end{itemize}
\end{mdframed}

% ══════════════════════════════════════════════════════════════
\section{Estimated Resources}
% ══════════════════════════════════════════════════════════════

\begin{center}
\begin{tabular}{p{5cm}p{8cm}}
\toprule
\textbf{Resource} & \textbf{Estimate} \\
\midrule
Antibody reagents &
FOXA1 CST D4E2 \#53528:
${\sim}$\$400/100 tests.
EZH2 CST D2H1 \#5246:
${\sim}$\$400/100 tests.
For 400 cases:
${\sim}$\$1,500--\$3,000 total. \\[0.3em]
IHC consumables &
${\sim}$\$2,000--\$4,000
(substantially reduced in TMA
format). \\[0.3em]
Pathologist scoring time &
${\sim}$40--60 hours whole section;
substantially reduced with TMA. \\[0.3em]
Statistical analysis &
${\sim}$10--20 hours. \\[0.3em]
\textbf{Total (consumables)} &
\textbf{\$5,000--\$15,000} \\[0.3em]
\textbf{Timeline} &
\textbf{3--6 months from tissue
access to submission} \\
\bottomrule
\end{tabular}
\end{center}

% ══════════════════════════════════════════════════════════════
\section{The Equity Argument}
% ══════════════════════════════════════════════════════════════

\begin{mdframed}[backgroundcolor=equitybox,
                 linecolor=orange!60!black,
                 linewidth=1pt]
Approximately 1.2--1.5 million breast cancer
patients per year currently receive no
molecular subtype information --- not because
the science does not exist, but because it
exists in a form requiring infrastructure
they do not have.

A validated FOXA1/EZH2 IHC protocol would
resolve the LumA/LumB ambiguity and the
TNBC/claudin-low ambiguity from two antibody
stains at \$50--\$100 per case, in any
laboratory that processes biopsies, on the
same day as the biopsy workup. No shipping.
No RNA. No waiting.

Your institution's participation in the
calibration study that enables this would
be directly credited in the first publication
of a diagnostic tool with this scale of
global health impact.

The calibration study costs \$5,000--\$15,000
in consumables. The potential reach on
validation is every breast cancer patient on
earth who receives a pathology workup.
\end{mdframed}

% ══════════════════════════════════════════════════════════════
\section{Published Evidence and Repository}
% ══════════════════════════════════════════════════════════════

\begin{center}
\begin{tabular}{p{4cm}p{9cm}}
\toprule
CS-LIT-1 (ratio axis) &
\href{https://doi.org/10.5281/zenodo.18883922}
{doi.org/10.5281/zenodo.18883922} \\
CS-LIT-22 (IHC tool) &
\href{https://doi.org/10.5281/zenodo.18892788}
{doi.org/10.5281/zenodo.18892788} \\
Full series (17 DOIs) &
\href{https://github.com/Eric-Robert-Lawson/attractor-oncology}
{github.com/Eric-Robert-Lawson/attractor-oncology} \\
\bottomrule
\end{tabular}
\end{center}

% ══════════════════════════════════════════════════════════════
\section{Next Steps}
% ═══════════════════════════════��══════════════════════════════

\begin{enumerate}
    \item \textbf{Contact:}
    \href{mailto:OrganismCore@proton.me}
    {OrganismCore@proton.me}
    \item \textbf{Confirm} approximate number
    of cases with PAM50 data and whether
    tissue is held in whole-section FFPE
    or TMA format
    \item \textbf{Review} attached Protocol
    Specification v3.0
    \item \textbf{Proceed} to institutional
    ethics review
\end{enumerate}

\vspace{0.5em}
\noindent\rule{\textwidth}{0.4pt}
\vspace{0.3em}

\begin{center}
\small
Eric Robert Lawson \textbar{} OrganismCore
\textbar{}
\href{mailto:OrganismCore@proton.me}
{OrganismCore@proton.me} \textbar{}
ORCID:
\href{https://orcid.org/0009-0002-0414-6544}
{0009-0002-0414-6544}\\
\href{https://github.com/Eric-Robert-Lawson/attractor-oncology}
{github.com/Eric-Robert-Lawson/attractor-oncology}
\textbar{}
CS-LIT-1:
\href{https://doi.org/10.5281/zenodo.18883922}
{10.5281/zenodo.18883922}
\textbar{}
CS-LIT-22:
\href{https://doi.org/10.5281/zenodo.18892788}
{10.5281/zenodo.18892788}
\end{center}

\end{document}
